\documentclass[10pt,a4paper]{amsart}
\usepackage{amsmath,amssymb,amsthm,graphicx}
% Retained minimal preamble for the Stacks Project chapter more-algebra.tex (More on Algebra).
% The chapter file pinned-original-source.tex carries only \input{preamble} on its first line, so no
% preamble text lives inside the chapter itself.  The Stacks Project uses one single shared preamble for
% every chapter; the snapshot stages it identically as sources/stacks-chow/preamble.tex and
% sources/stacks-etale-cohomology/preamble.tex.  The block below is copied mechanically, line for line,
% from sources/stacks-chow/preamble.tex lines 14-18 (the Xy-pic setup) and 154-217 (the native theorem
% declarations together with the mathematical macros).  Nothing was added, removed or redefined; the
% wrapper class and the amsmath/amssymb/amsthm/graphicx base packages are supplied by the assembler.
% Chapter: more-algebra.tex -- More on Algebra.
% For commutative diagrams we use Xy-pic
\usepackage[all]{xy}

% We use 2cell for 2-commutative diagrams.
\xyoption{2cell}

% Theorem environments.
%
\theoremstyle{plain}
\newtheorem{theorem}[subsection]{Theorem}
\newtheorem{proposition}[subsection]{Proposition}
\newtheorem{lemma}[subsection]{Lemma}

\theoremstyle{definition}
\newtheorem{definition}[subsection]{Definition}
\newtheorem{example}[subsection]{Example}
\newtheorem{exercise}[subsection]{Exercise}
\newtheorem{situation}[subsection]{Situation}

\theoremstyle{remark}
\newtheorem{remark}[subsection]{Remark}
\newtheorem{remarks}[subsection]{Remarks}

\numberwithin{equation}{subsection}

% Macros
%
\def\lim{\mathop{\mathrm{lim}}\nolimits}
\def\colim{\mathop{\mathrm{colim}}\nolimits}
\def\Spec{\mathop{\mathrm{Spec}}}
\def\Hom{\mathop{\mathrm{Hom}}\nolimits}
\def\Ext{\mathop{\mathrm{Ext}}\nolimits}
\def\SheafHom{\mathop{\mathcal{H}\!\mathit{om}}\nolimits}
\def\SheafExt{\mathop{\mathcal{E}\!\mathit{xt}}\nolimits}
\def\Sch{\mathit{Sch}}
\def\Mor{\mathop{\mathrm{Mor}}\nolimits}
\def\Ob{\mathop{\mathrm{Ob}}\nolimits}
\def\Sh{\mathop{\mathit{Sh}}\nolimits}
\def\NL{\mathop{N\!L}\nolimits}
\def\CH{\mathop{\mathrm{CH}}\nolimits}
\def\proetale{{pro\text{-}\acute{e}tale}}
\def\etale{{\acute{e}tale}}
\def\QCoh{\mathit{QCoh}}
\def\Ker{\mathop{\mathrm{Ker}}}
\def\Im{\mathop{\mathrm{Im}}}
\def\Coker{\mathop{\mathrm{Coker}}}
\def\Coim{\mathop{\mathrm{Coim}}}

% Boxtimes
%
\DeclareMathSymbol{\boxtimes}{\mathbin}{AMSa}{"02}

%
% Macros for moduli stacks/spaces
%
\def\QCohstack{\mathcal{QC}\!\mathit{oh}}
\def\Cohstack{\mathcal{C}\!\mathit{oh}}
\def\Spacesstack{\mathcal{S}\!\mathit{paces}}
\def\Quotfunctor{\mathrm{Quot}}
\def\Hilbfunctor{\mathrm{Hilb}}
\def\Curvesstack{\mathcal{C}\!\mathit{urves}}
\def\Polarizedstack{\mathcal{P}\!\mathit{olarized}}
\def\Complexesstack{\mathcal{C}\!\mathit{omplexes}}
% \Pic is the operator that assigns to X its picard group, usage \Pic(X)
% \Picardstack_{X/B} denotes the Picard stack of X over B
% \Picardfunctor_{X/B} denotes the Picard functor of X over B
\def\Pic{\mathop{\mathrm{Pic}}\nolimits}
\def\Picardstack{\mathcal{P}\!\mathit{ic}}
\def\Picardfunctor{\mathrm{Pic}}
\def\Deformationcategory{\mathcal{D}\!\mathit{ef}}

% --- appended 2026-10-12 (Chief): the project own definitions of the reference/slogan/history
% display environments, copied mechanically from the shared Stacks preamble lines 8-12
% (\usepackage{verbatim} plus the three \newenvironment lines, which the chapter preamble keeps
% verbatim from the project preamble). Nothing added, removed or redefined: these environments
% carry no mathematical content and the project itself typesets them as comments. Without these
% three lines a counted statement whose own text uses them fails with Environment slogan undefined.
\usepackage{verbatim}
\newenvironment{reference}{\comment}{\endcomment}
\newenvironment{slogan}{\comment}{\endcomment}
\newenvironment{history}{\comment}{\endcomment}
\usepackage{cleveref}
\crefname{originallemma}{original lemma}{original lemmas}
\Crefname{originallemma}{Original lemma}{Original lemmas}
\makeatletter
\newlabel{derived-lemma-compute-ext-resolutions}{{06XR}{source}{}{sourcepointer-1}{}}
\newlabel{derived-lemma-compute-ext-resolutions@cref}{{[originallemma][0][]06XR}{[1][0][]source}{}{}{}}
\makeatother
\crefname{originallemma}{original lemma}{original lemmas}
\Crefname{originallemma}{Original lemma}{Original lemmas}
\makeatletter
\newlabel{derived-lemma-negative-exts}{{06XS}{source}{}{sourcepointer-2}{}}
\newlabel{derived-lemma-negative-exts@cref}{{[originallemma][0][]06XS}{[1][0][]source}{}{}{}}
\makeatother
\crefname{originallemma}{original lemma}{original lemmas}
\Crefname{originallemma}{Original lemma}{Original lemmas}
\makeatletter
\newlabel{derived-remark-truncation-distinguished-triangle}{{08J5}{source}{}{sourcepointer-3}{}}
\newlabel{derived-remark-truncation-distinguished-triangle@cref}{{[originallemma][0][]08J5}{[1][0][]source}{}{}{}}
\makeatother
\crefname{originallemma}{original lemma}{original lemmas}
\Crefname{originallemma}{Original lemma}{Original lemmas}
\makeatletter
\newlabel{lemma-projective-amplitude}{{17414}{source}{}{sourcepointer-4}{}}
\newlabel{lemma-projective-amplitude@cref}{{[originallemma][0][]17414}{[1][0][]source}{}{}{}}
\makeatother
\crefname{originallemma}{original lemma}{original lemmas}
\Crefname{originallemma}{Original lemma}{Original lemmas}
\makeatletter
\newlabel{lemma-splitting-unique}{{19704}{source}{}{sourcepointer-5}{}}
\newlabel{lemma-splitting-unique@cref}{{[originallemma][0][]19704}{[1][0][]source}{}{}{}}
\makeatother
\title{Injectivity of an Ext group along injective module maps forces a splitting of the truncation triangle}
\author{The Stacks Project Authors}
\date{stacks-project:more-algebra.tex:local-snapshot-20261012; GNU FDL 1.2 or later}
\begin{document}
\maketitle
\setcounter{subsection}{1}
\noindent\textit{Editorial note.} All delivered mathematical text is the author's own, copied from the pinned Stacks Project chapter file more-algebra.tex; no statement or proof step was rewritten, reordered or removed.\par\medskip
\noindent\textit{Editorial note.} The chapter file pinned-original-source.tex carries only \textbackslash{}input\{preamble\} on its first line, so no preamble text lives inside the chapter. The retained file sources/stacks-more-algebra/retained-preamble.tex is a mechanical, line-for-line copy of the project's single shared preamble as staged in the snapshot; nothing was added, removed or redefined and no package was added.\par\medskip
\noindent\textit{Editorial note.} The source URL is the source-level Stacks Project URL, as required when more than one card is produced from one source; the per-claim tag is carried by claim\_key.\par\medskip
\noindent\textit{Editorial note.} The author writes '(this is where we use that K is bounded above)' inside the counted proof span, reproduced verbatim. The five references inside the counted proof span are recorded as non-load-bearing navigation pointers.\par\medskip
\noindent\textit{Allowed standard prerequisites (editorial).}\begin{itemize}
\item Truncation distinguished triangles and the computation of Ext groups of a truncation. invoked at proof lines 19929-19939. Supplies the exact sequence whose last term is Hom\_R(H\textasciicircum{}a(K), M).
\item Vanishing of Ext of a bounded above complex out of an injective module, and the fact that every module embeds in an injective one. used at proof lines 19948-19952. Is what turns the hypothesis about injective maps into vanishing for all modules.
\item The criterion that vanishing of Ext\textasciicircum{}\{-a\} against all modules is equivalent to projective amplitude in [a+1, b], and the unique-splitting lemma. invoked at proof lines 19953-19957. Produces the amplitude bound and then the unique direct sum decomposition.
\end{itemize}\medskip
\ifdefined\hypertarget\hypertarget{sourcepointer-1}{}\fi
\noindent\textit{Source reference.} Original lemma 06XR: focused-seeds/AGP/sources/stacks-snapshot-20261011/tags.txt, source lines 8860--8860. Cross-chapter navigation pointer into the Stacks Project; the standard fact it supplies is applied inside the counted proof span and the argument resting on it stays there.\par\medskip
\ifdefined\hypertarget\hypertarget{sourcepointer-2}{}\fi
\noindent\textit{Source reference.} Original lemma 06XS: focused-seeds/AGP/sources/stacks-snapshot-20261011/tags.txt, source lines 8861--8861. Cross-chapter navigation pointer into the Stacks Project; the standard fact it supplies is applied inside the counted proof span and the argument resting on it stays there.\par\medskip
\ifdefined\hypertarget\hypertarget{sourcepointer-3}{}\fi
\noindent\textit{Source reference.} Original lemma 08J5: focused-seeds/AGP/sources/stacks-snapshot-20261011/tags.txt, source lines 10838--10838. Cross-chapter navigation pointer into the Stacks Project; the standard fact it supplies is applied inside the counted proof span and the argument resting on it stays there.\par\medskip
\ifdefined\hypertarget\hypertarget{sourcepointer-4}{}\fi
\noindent\textit{Source reference.} Original lemma 17414: sources/stacks-more-algebra/pinned-original-source.tex, source lines 17414--17414. Same-chapter navigation pointer; the statement it names is standard for the section and the argument that consumes it stays inside the counted proof span.\par\medskip
\ifdefined\hypertarget\hypertarget{sourcepointer-5}{}\fi
\noindent\textit{Source reference.} Original lemma 19704: sources/stacks-more-algebra/pinned-original-source.tex, source lines 19704--19704. Same-chapter navigation pointer; the statement it names is standard for the section and the argument that consumes it stays inside the counted proof span.\par\medskip
\begin{lemma}
\label{lemma-split-using-ext-injective}
Let $R$ be a ring. Let $K \in D^-(R)$. Let $a \in \mathbf{Z}$.
Assume that for any injective $R$-module map $M \to M'$ the map
$\Ext^{-a}_R(K, M) \to \Ext^{-a}_R(K, M')$ is injective.
Then there is a unique direct sum decomposition
$K \cong \tau_{\leq a}K \oplus \tau_{\geq a + 1}K$
and $\tau_{\geq a + 1}K$ has projective-amplitude in $[a + 1, b]$
for some $b$.
\end{lemma}

\begin{proof}
Consider the distinguished triangle
$$
\tau_{\leq a}K \to K \to \tau_{\geq a + 1}K \to (\tau_{\leq a}K)[1]
$$
in $D(R)$, see Derived Categories, Remark
\ref{derived-remark-truncation-distinguished-triangle}.
Observe that
$\Ext^{-a}_R(\tau_{\leq a}K, M) = \Hom_R(H^a(K), M)$
and
$\Ext^{-a - 1}_R(\tau_{\leq a}K, M) = 0$, see
Derived Categories, Lemma \ref{derived-lemma-negative-exts}. Thus
the long exact sequence of $\Ext$ gives an exact sequence
$$
0 \to
\Ext^{-a}_R(\tau_{\geq a + 1}K, M) \to
\Ext^{-a}_R(K, M) \to
\Hom_R(H^a(K), M)
$$
functorial in the $R$-module $M$.
Now if $I$ is an injective $R$-module, then
$\Ext^{-a}_R(\tau_{\geq a + 1}K, I) = 0$ for example by
Derived Categories, Lemma \ref{derived-lemma-compute-ext-resolutions}.
Since every module injects into an injective module,
we conclude that $\Ext^{-a}_R(\tau_{\geq a + 1}K, M) = 0$
for every $R$-module $M$. By Lemma \ref{lemma-projective-amplitude}
we conclude that $\tau_{\geq a + 1}K$ has projective-amplitude
in $[a + 1, b]$ for some $b$ (this is where we use that $K$
is bounded above). We obtain the splitting by
Lemma \ref{lemma-splitting-unique}.
\end{proof}

\end{document}
